\documentclass{ieeeaccess}
\usepackage{cite}
\usepackage{amsmath,amssymb,amsfonts}
\usepackage{algorithm}
\usepackage{algorithmicx}
\usepackage{algpseudocode}
\usepackage{graphicx}
\usepackage{textcomp}
\usepackage{booktabs}
\usepackage{multirow}
\usepackage{tabularx}
\usepackage{float}
\usepackage{xcolor}
\usepackage{soul}
\sethlcolor{yellow}
\def\BibTeX{{\rm B\kern-.05em{\sc i\kern-.025em b}\kern-.08em
    T\kern-.1667em\lower.7ex\hbox{E}\kern-.125emX}}

\begin{document}
\history{Date of publication xxxx 00, 0000, date of current version xxxx 00, 0000.}
\doi{10.1109/ACCESS.2026.0000000}

\title{Managing In-Memory Entity Blocks in a Streaming Context}
\author{\uppercase{Igor de Sousa Pereira}\authorrefmark{1}, \uppercase{Carlos Eduardo S. Pires}\authorrefmark{1}, and \uppercase{Tiago Brasileiro Ara\'ujo}\authorrefmark{2}\authorrefmark{1}}

\address[1]{Department of Systems and Computing, Federal University of Campina Grande, Para\'iba, Brazil (e-mail: igor.pereira@copin.ufcg.edu.br; cesp@dsc.ufcg.edu.br)}
\address[2]{Federal Institute of Para\'iba, Soledade, Para\'iba, Brazil (e-mail: tiago.brasileiro@ifpb.edu.br)}

\markboth
{Pereira \headeretal: Managing In-Memory Entity Blocks in a Streaming Context}
{Pereira \headeretal: Managing In-Memory Entity Blocks in a Streaming Context}

\corresp{Corresponding author: Igor de Sousa Pereira (e-mail: igor.pereira@copin.ufcg.edu.br).}

\begin{abstract}
With the growth of online services on the Internet, a large volume of data has been generated, making it difficult to identify records that refer to the same real-world entity, especially in continuous data stream contexts. As a result, Entity Resolution (ER) techniques have had to adapt to the streaming scenario by using entity blocks in order to eliminate unnecessary comparisons, facilitate the identification of duplicate entities across multiple data sources, and produce sets of potentially similar records. \hl{This paper evaluates three eviction policies for managing in-memory entity blocks in streaming ER: per-block time-aware entity removal, full block clearing at a capacity trigger, and a hybrid policy combining global age-based removal with local block clearing.} In addition, the proposed techniques are evaluated in noisy ER scenarios, where spelling and typographical variations may affect blocking quality. \hl{The evaluation reports blocking time, matching time, candidate-pair counts, Recall, and Precision relative to the Approximate Blocking baseline; conclusions about matching quality are restricted to validated metric calculations.}
\end{abstract}

\begin{keywords}
Blocking, Entity Resolution, Hash, Resource Management, Streaming.
\end{keywords}

\titlepgskip=-15pt
\maketitle
\section{Introduction} \label{sec:introducao}

With the growing volume of data generated today, several applications have been developed with the capability to store information originating from diverse sources, such as private communication services, including electronic messages, and web-based commercial platforms, such as digital newspapers and e-commerce systems. This continuous increase in data is also driven by the emergence and widespread adoption of social networks, such as Facebook\footnote{https://www.facebook.com/}, X\footnote{https://x.com/}, and LinkedIn\footnote{https://br.linkedin.com/}, which facilitate user interaction and generate large volumes of data arising from interactions similar to those occurring in the physical world \cite{10.1111/j.1083-6101.2007.00393.x}.

In this context, Entity Resolution (ER) tasks \cite{christen2012} have been developed with the goal of identifying duplicate information or records referring to the same real-world entity across distinct data sources and heterogeneous domains. ER has been widely applied in areas such as Healthcare \cite{10.1145/3448016.3457238}, to identify symptoms and their corresponding diseases, and Social Networks \cite{doi:10.1504/IJWGS.2017.085167}, to group mentions of the same real-world event originating from multiple platforms, such as \textit{Twitter} (currently X) and \textit{SogouLab}\footnote{https://www.sogou.com/?fr=labs}.

In recent years, ER tasks have increasingly been deployed in data streaming environments to address the challenges of continuously generated data \cite{karapiperis2023, Gazzarri2021, DBLP:conf/edbt/GazzarriH23}. In this context, ER poses specific challenges: (i) operate under memory constraints and continuous processing requirements, (ii) make incremental and often irreversible decisions as new records arrive, and (iii) copy with temporal evolution and the arrival of late or inconsistent information. 

To illustrate, consider a stream of scientific article records continuously submitted by multiple research institutions and academic platforms, such as Google Scholar, Scopus, and DBLP. Each record contains attributes including an identifier, title, authors, publication venue, and publication year. For example, one record identified as ``0HMk'' corresponds to the article \textit{Too Much Middleware}, authored by M.~Stonebraker and published in \textit{SIGMOD Record} in 2002. Another record, identified as ``ixKf'', corresponds to the article \textit{Initiation of crazes in polystyrene}, authored by A.~S.~Argon and J.~G.~Hannoosh and published in \textit{Phil.\ Mag.} in 1975.

The challenge consists of identifying potential matches between records that may refer to the same article or the same author, even in the presence of variations in writing styles, including name abbreviations (“J. Hannoosh” vs. “JG Hannoosh”), differences in title formatting (“Too Much Middleware” vs. “too much middleware”), or inconsistencies in source metadata (“2002” vs. “2020”).

In a streaming context, the ER task must operate in (near) real time, as new publications are continuously incorporated into the datasets \cite{araujo2026x}. This requires the ER task to be performed both efficiently and accurately in order to detect duplicate articles or authors and to notify a central system whenever record overlap occurs; for example, when the same article is indexed under different identifiers.

Typically, the ER task is carried out in three main stages \cite{christen2012}. The first stage, referred to as \textit{Data Preprocessing}, is responsible for cleaning and standardizing information originating from data sources, thereby facilitating subsequent processing steps. The second stage, known as \textit{Blocking}, groups potentially similar entities into blocks based on shared attribute values, referred to as blocking keys. Such keys may correspond, for instance, to the initial characters of author names or to the publication year of an article.

The third stage, referred to as \textit{Matching}, aims to determine the degree of similarity between pairs of entities that belong to the same block. Comparisons are performed based on the attribute values of the entities involved. To determine whether two entities are similar, different similarity or distance functions can be applied to attribute values, such as Jaccard similarity and Euclidean distance. Figure~\ref{fig:fluxo_re} illustrates the overall workflow of the ER task. In particular, this work focuses on the Blocking stage, which is especially important in streaming environments due to the need for efficient and accurate real-time processing.

\begin{figure}[h]
    \centering
    \includegraphics[width=0.8\linewidth]{images/fluxo_re-2.png}
    \caption{Basic ER workflow.}
    \label{fig:fluxo_re}
\end{figure}

Beyond this, data originating from heterogeneous sources exhibit structural and semantic variations that directly impact the blocking stage, which is considered a critical step of the ER task \cite{ARAUJO2025102506}. Inadequate blocking may produce excessively large blocks, increasing the number of redundant comparisons and impairing the efficiency of the matching process. For this reason, techniques based on hash functions are widely adopted \cite{10.1145/3297280.3299730}, as they enable the grouping of potentially similar records without the need for exhaustive comparisons, even in the presence of noise or minor textual variations. Among these techniques, Locality-Sensitive Hashing (LSH) stands out \cite{lshbased01, 10.5555/645925.671516}, as it increases the probability of mapping similar entities to the same blocks, thereby reducing the search space.

Due to the characteristics of streaming environments,  such as the high arrival rate of entities and the often unbounded data volume, it becomes infeasible to compare all pairs of entities accumulated over time \cite{10.1145/3448016.3457238}. As shown in \ref{exp6}, when no eviction strategy is applied, the matching processing time in the MusicBrainz dataset\footnote{https://dbs.uni-leipzig.de/research/projects/benchmark-datasets-for-entity-resolution} increases by approximately 24 times, reaching nearly 8 hours, compared to only 20 minutes when using a simple eviction strategy.

To address this issue, approaches such as those proposed in \cite{Gazzarri2021,10.1145/3341105.3375776} aim to manage entity blocks efficiently by limiting the number of blocks kept in memory and applying cleaning strategies to remove older blocks. These approaches seek to mitigate the challenges imposed by the massive scale and continuous flow of data in streaming environments. However, such techniques are generally based on empirical block management methods, which hinder formal validation and may compromise the accuracy of the matching stage.

An alternative to this problem is the use of an in-memory block structure, such as the one presented in \cite{10.1109/TKDE.2024.3361022}, called AB, which stores information such as the maximum number of entities per block. However, existing techniques exhibit several limitations: the entity eviction process is based on simple removals and does not take into account the lifetime of entities within blocks, nor does it adequately address the problem of saturated blocks, i.e., blocks that accumulate a large number of irrelevant entities and may compromise the efficiency of the ER task. This simplistic approach can lead to memory waste and increased computational cost during entity comparisons.

In this paper, we propose eviction and block management strategies for streaming ER settings aimed at addressing the inefficiencies caused by simplistic eviction policies, saturated blocks, and excessive comparison costs. \hl{The evaluated policies use entity age and block capacity as eviction criteria; they do not estimate semantic relevance or guarantee lower memory use or better matching quality.} In this context, the research question addressed by this paper is: “\hl{How do the evaluated eviction strategies affect blocking and matching execution time, candidate-pair counts, Recall, and Precision in the examined streaming ER settings?}”

The main contributions of this work focus on the development and evaluation of three complementary eviction strategies aimed at the efficient maintenance of blocking structures for streaming ER. \hl{The strategies are evaluated using processing time, candidate-pair counts, Recall, and Precision; they are not claimed to preserve matching quality beyond what is supported by validated measurements.} In summary, this work presents:
 
\begin{itemize}\small
    \item \textbf{Entity Eviction:} a block-dependent eviction strategy that removes the oldest entities when a block reaches its maximum capacity, prioritizing more recent and potentially more relevant records (Section~\ref{subsec:descarte-entidades});
    
    \item \textbf{Block Eviction:} a strategy that removes entire blocks when they exceed a predefined size threshold, preventing the formation of poorly discriminative structures with low matching potential (Section~\ref{subsec:descarte-blocos});
    
    \item \textbf{Global Entity Eviction:} \hl{a hybrid strategy that combines global time-based removal of entities outside the retention window with local block clearing when a block reaches its capacity limit} (Section~\ref{subsec:descarte-entidades-global}).
    
\end{itemize}

The remainder of this paper is organized as follows: Section~\ref{sec:background} presents the fundamental concepts required to understand the proposed solution; Section~\ref{sec:trabRelacionados} discusses related work; Section~\ref{sec:estrategias-descarte-entidades} introduces the proposed strategies; Section~\ref{sec:avaliacao} presents the experimental evaluation and a discussion of the results. Finally, Section~\ref{sec:conclusao} concludes the paper and outlines directions for future work.

\section{Background} \label{sec:background}

This section introduces the fundamental concepts underlying this work. We first review the blocking process in ER, which reduces the number of comparisons by grouping potentially similar records. We then discuss ER in streaming environments, highlighting the challenges of continuous data processing. Finally, we present Approximate Blocking (AB), the technique adopted in this study to enable incremental blocking over data streams.

\subsection{Blocking in Entity Resolution}
The Blocking stage plays a fundamental role in the ER task, as it acts as an initial filter that reduces the number of entity comparisons. Before any matching technique is applied, it is necessary to identify subsets of entities that are likely to represent the same real-world object \cite{christen2012}. To this end, blocks are created to restrict comparisons only to pairs of entities belonging to the same block, eliminating unnecessary checks between clearly distinct records. In this way, blocking groups potentially similar entities and enables the system to operate in a scalable manner even in the presence of large data volumes.

There are two main types of blocking: exact and approximate. Exact blocking forms blocks based on identical keys and is suitable when attribute values are stable and free of noise. AB, on the other hand, is designed to handle noisy data, i.e., records whose attributes may contain errors, inconsistencies, or formatting variations. As a result, values such as ``Catherine'' and ``Katherine,'' for example, may be assigned to the same block, increasing coverage and reducing the likelihood of missing similar entities. At the end of this stage, a reduced set of candidate entity pairs is obtained to be compared in the Matching stage.

To deal with textual variations, it is common to employ hash functions \cite{10.1145/3297280.3299730}, which convert attribute values of records into compact numerical codes. These codes allow similar records to be grouped without requiring direct comparisons among all entities in the dataset. Traditional hashes generate deterministic values and are suitable for exact blocking. In scenarios where attribute values contain textual variations, noise, or inconsistencies, \textit{Locality-Sensitive Hashing} (LSH) is commonly employed. LSH is a family of \textit{hash} functions designed to preserve similarity between vectors representing records, typically constructed from textual attributes. In this case, records with similar values have a higher probability of collision and may be assigned to the same block. This technique reduces the number of comparisons performed during the Matching stage, which is particularly advantageous in scenarios characterized by high data volume and dynamism, such as continuous data streams \cite{10.1145/997817.997857}.

In the streaming context, the challenges become even more pronounced. Unlike static scenarios, in which the data domain is known in advance and structures can be prepared beforehand, streaming data arrive continuously, dynamically, and potentially infinitely \cite{araujo2026x}. This imposes constraints on time, memory, and update capacity. In such environments, the ER process must be incremental, where each new entity is processed without the possibility of re-executing the entire blocking and matching pipeline, which would be computationally infeasible \cite{DONASCIMENTO2018335, DBLP:journals/corr/RenC17}.

Moreover, real-world streams may involve data originating from distinct sources with heterogeneous structures. Works such as \cite{10.14778/2856318.2856326, Papadakis2016ScalingER} investigate blocking strategies in such scenarios, comparing two main approaches: \textit{schema-aware} and \textit{schema-agnostic}. \textit{schema-aware} approaches exploit structural information, such as attribute names or types, to generate more interpretable and specialized blocking keys. When data sources share a common schema, \textit{schema-aware} approaches can achieve competitive performance in terms of reducing the number of comparisons. However, experimental results indicate that their performance tends to be unstable, particularly with respect to \textit{Recall}, and degrades significantly in heterogeneous contexts, where schemas diverge or are unknown \cite{10.14778/2856318.2856326}.

In contrast, \textit{schema-agnostic} approaches do not require prior knowledge of the schema, automatically extracting tokens from the textual content of attributes to form blocking signatures. This flexibility makes such strategies scalable and suitable for distributed environments or highly heterogeneous streams. Nevertheless, they may generate a larger number of comparisons and reduce precision when the extracted tokens do not adequately capture the semantics of the entities.

Therefore, the choice of a blocking technique (exact or approximate, \textit{schema-aware} or \textit{schema-agnostic}) depends on factors such as the presence of noise, structural variability of the schema, data volume, and the need for continuous updates. In streaming environments, robust and flexible techniques become essential to preserve the quality of ER in the presence of dynamic, heterogeneous, and potentially noisy data.

\subsection{Entity Resolution with Streaming Data Sources}

Studies have shown that, in 2020, the amount of data generated daily reached approximately 44 zettabytes, which is equivalent to about 44 billion terabytes \cite{Berisha2022}. As this volume continues to grow, storage systems must cope with such massive amounts of information, requiring fast and efficient data structures to ingest and manage these data.

The importance of data streaming is evidenced by its adoption across several domains, such as Healthcare, Finance, E-commerce, Social networks, and the Internet of Things (IoT) \cite{DHARAVATH201593, gentile2017entity, 10.1145/3448016.3457238, doi:10.1504/IJWGS.2017.085167, Sergey2012JointLU}. An illustrative example can be found in environmental monitoring through continuous streams of satellite imagery. In such scenarios, satellites transmit data in near real time, which are instantaneously processed to detect patterns and changes on the Earth’s surface. This enables the monitoring of deforestation, urban expansion, extreme weather events, and even agricultural indicators. Rather than relying on periodic analyses, streaming data allows the generation of immediate insights and supports timely decision-making.

When the ER task is applied in a streaming data context, several challenges arise, including: (i) maintaining entity blocks under continuous data arrival; (ii) managing memory efficiently under storage constraints; (iii) ensuring low latency in continuous processing;  (iv) adapting to variations in data volume and attribute characteristics; and (v) balancing accuracy and performance given the large number of comparisons.

Unlike traditional ER, it is necessary to optimize the process to handle a continuous and potentially infinite data stream, requiring (near) real-time processing. One of the main challenges lies in the efficient management of entity blocks in memory due to storage limitations, which demands strict control mechanisms to remove obsolete blocks without compromising the quality of matching results.

These challenges manifest throughout the entire ER pipeline. In the blocking stage, the creation and maintenance of blocks become complex due to the continuous data flow and the need for frequent updates under memory constraints. In the matching stage, data unpredictability, expressed through variations in volume, temporal distribution, and entity attributes, requires efficient and scalable algorithms capable of performing a large number of comparisons within tight time constraints. In this context, it is crucial to balance computational performance and matching quality, ensuring that relevant pairs are identified without incurring excessive processing costs.

Among these challenges, the present work focuses on the implementation of optimized blocking techniques for streaming scenarios, enabling the dynamic creation and maintenance of entity blocks that automatically adapt to variations in data volume and characteristics.

\subsection{Approximate Blocking}

The \textit{Approximate Blocking} (AB) technique \cite{10.1109/TKDE.2024.3361022} proposes a set of algorithms based on randomized mechanisms for block generation in streaming ER, focusing on computational efficiency and constrained memory usage. The \textit{AB} technique adopts block maintenance and replacement policies to ensure prioritized retention of the most relevant entity blocks, while integrating AB to provide robustness against noise and data variations.

Despite its benefits, the use of randomization in block generation may affect the consistency of the candidate pair set produced by the blocking stage, thereby impacting the final quality of ER. AB is characterized by relaxing exact matching criteria through similarity metrics, increasing recall (coverage) in noisy data scenarios at the cost of generating a larger number of false positives. The higher number of false positives arises because the blocking criterion is no longer strictly deterministic, allowing non-matching records with sufficient superficial similarity to be grouped into the same block. This family of techniques advances the feasibility of ER over high-velocity and high-volume data streams, representing an important contribution to real-time data integration systems.

Figure \ref{fig:fluxo_aprxmt} illustrates an example of an ER architecture based on AB \cite{10.1109/TKDE.2024.3361022}, highlighting the processing flow from the ingestion of heterogeneous records to the identification of entity matches.In this scenario, differences between attribute sets, such as $\{a_1, a_2\}$ and $\{a_3, a_4\}$, indicate variations across distinct data sources. The example is therefore characterized as \textit{schema-agnostic}, since it does not rely on a common attribute schema across data sources, operating directly on textual representations of records. In addition, the architecture is incremental, and capable of handling frequent data ingestion.

\begin{figure}[h]
  \centering
  \includegraphics[width=0.9\linewidth]{images/flx_apxmt-1.png}
  \caption{Approximate Blocking-based ER Workflow.}
  \label{fig:fluxo_aprxmt}
\end{figure}

Based on this general architecture, the operational workflow proceeds as follows: two datasets are provided to the process. The first dataset is used in the blocking stage, while the second is intended for the \textit{matching} stage. Initially, Dataset 1 is read and subjected to a preprocessing step, which is responsible for reading and organizing entities in main memory. Next, the hash generation stage is initiated, whose objective is to map entities to blocking keys efficiently, reducing the comparison space in the subsequent stage. At this step, for each entity, \(R\) distinct hash values are generated; in this example, \(R = 4\). Each hash value is used exclusively as a blocking key, allowing entities that share the same value to be allocated to the same block. As a result, a single entity, such as \(e_1\), may be associated with multiple blocks (\(b_1\), \(b_2\), \(b_3\), \(b_7\)), increasing the likelihood that potentially similar entities share at least one block. It is worth noting that such collisions are approximate in nature and do not imply a direct correspondence between entities; rather, they serve to define candidate sets that will later be evaluated during the comparison stage.

In the second stage, corresponding to \textit{matching}, Dataset 2 is read, pre-processed, and subjected to the same hash generation process, although with a smaller number of hash values generated per entity compared to the blocking stage. In the example of Figure~\ref{fig:fluxo_aprxmt}, entity $e_4$ is introduced, generating two hash values. The \textit{matching} process then verifies whether these hashes correspond to existing blocking keys. If key repetitions occur across entities, that is, if the hashes of $e_4$ coincide with blocks containing $e_1$ and this overlap exceeds a threshold defined by $\alpha$, a \textit{match} between the analyzed entities is determined. The pseudo-code illustrating the execution logic of \textit{AB} is provided in Algorithm~\ref{alg:approx_blocking}.

The main idea of this technique is to use the co-occurrence of values to determine the similarity between entities. In the AB technique, co-occurrence represents the frequency with which pairs of similar records appear together across different approximate blocks. This characteristic enables the identification of true matches even in the presence of textual variations or noise in the data, overcoming the limitations of exact blocking. By combining similarity metrics with randomized structures, the technique balances accuracy and efficiency, ensuring broad coverage of relevant candidate pairs. Consequently, value co-occurrence makes the blocking process more robust and better suited to continuous and noisy data scenarios, such as \textit{streaming}.

\begin{algorithm}[h]
\caption{Resolving a query record $q$ using AB}
\label{alg:approx_blocking}
\begin{algorithmic}[1]

\Require Array $T$ of inverted indices; Number of sampled indices $R'$; Threshold $\theta$.

\State $matchingpairs \gets \emptyset$
\State $collisionCount \gets \emptyset$ \Comment{Collision count map}

\For{$i = 1$ to $R'$}
    \State $j \gets \text{Random}(1, R)$
    \State $k \gets h_j(q)$ \Comment{Apply MinHash with seed $j$}
    \State $Ids \gets T[j][k]$ \Comment{Retrieve IDs in block $k$}

    \For{each $Id \in Ids$}
        \If{$Id \notin matchingpairs$}
            \State $collisionCount[Id] \gets collisionCount[Id] + 1$
        \EndIf
        \State $c \gets collisionCount[Id]$
        \If{$\frac{c}{R'} \ge \theta$}
            \State $matchingpairs \gets matchingpairs \cup \{Id\}$
        \EndIf
    \EndFor
\EndFor

\State \Return $mp$

\end{algorithmic}
\end{algorithm}

The results reported in \cite{10.1109/TKDE.2024.3361022} indicate that using a small number of \textit{hash} functions is sufficient to maintain high levels of accuracy in entity matching. The results indicate that, even with a limited set of randomly selected indices, the technique preserves the quality of the associations. Although this reduction introduces a small additional error, such error is controllable and proportional to the number of sampled indices. Thus, the technique achieves a balance between computational efficiency and accuracy, ensuring consistent results even in scenarios with large data volumes.

Considering the hash generation process, let the set of entities be denoted by \( \mathcal{E} = \{ e_1, e_2, \dots, e_N \} \), where each entity is described by a unique identifier and an attribute vector \( e_i = (id_i, \mathbf{a}_i) \), with \( \mathbf{a}_i = (a_{i1}, a_{i2}, \dots, a_{iD}) \) representing the set of \( D \) observable attributes. During the hash generation phase, the goal is to obtain compact representations that preserve similarity relationships among entities in the original attribute space, reducing dimensionality and facilitating the identification of potential matches.

To this end, a family of \textit{hash} functions \( \mathcal{H} = \{ h_1, h_2, \dots, h_R \} \) is defined, where each function \(h_r : \mathcal{A}^D \rightarrow \mathcal{K}\) maps the attribute vector \( \mathbf{a}_i \) to a discrete value representing a similarity group identifier. These groups act as partitions of the data space, organizing entities such that those with identical or close \textit{hash} values are considered potentially similar. As a result, the ER process becomes more efficient, since comparisons are restricted to reduced subsets of entities where the likelihood of true similarity is higher.

By applying each function in the set of \textit{hash} functions to an entity, a projection vector is obtained:
\[
\mathbf{h}(e_i) = (h_1(\mathbf{a}_i), h_2(\mathbf{a}_i), \dots, h_R(\mathbf{a}_i))
\]

Each function \( h_r \) captures a distinct aspect of similarity between entities, understood as a different probabilistic projection over the set of \(q\)-grams that compose their textual representations. For example, one \textit{hash} function may emphasize matches associated with specific subsets of \(q\)-grams, while another may favor different combinations of these same elements. In this way, the joint use of multiple functions reduces the probability that similar pairs are separated due to noise or attribute variations. Consequently, employing a family of \textit{hash} functions enhances the technique’s ability to identify potential matches by producing multiple complementary groupings in the blocking space.

Accordingly, for each \textit{hash} function \( h_r \in \mathcal{H} \), an independent set of blocks \(\mathcal{B}_r = \{ b_{r,1}, b_{r,2}, \dots, b_{r,K_r} \}\)
is maintained, which stores the entities whose \textit{hash} values were generated specifically by function \( h_r \). Each block \( b_{r,k} \) groups the entities that, when processed by \( h_r \), were assigned to the same grouping index, such that:
\[
    b_{r,k} =\{e_i \in \mathcal{E}\mid h_r(\mathbf{a}_i)=k \}.  
\]

In other words, each function \( h_r \) defines its own grouping structure, responsible for clustering all entities processed under the same similarity projection. Thus, although different functions within the family may produce numerically similar \textit{hash} values, each one maintains its results in a distinct set of blocks, ensuring separation among the partitions defined by each hash projection.

This organization guarantees that, during processing, the \textit{hash} values produced by each \( h_r \) are stored in the corresponding position of the inverted index associated with that function, which integrates the global blocking structure composed of a collection of inverted indices. Such an organization enables consistent grouping of entities within each partition, facilitating efficient candidate retrieval during the matching stage. Consequently, the adopted blocking technique (based on a family of \textit{hash} functions) preserves cohesion among related entities and optimizes search and comparison operations, while also simplifying the management of multiple similarity partitions generated by different functions in the set \( \mathcal{H} \).

\begin{table}[!t]
\centering
\caption{\hl{Notation used in the blocking, matching, and eviction processes.}}
\label{tab:notation}
\scriptsize
\renewcommand{\arraystretch}{1.05}
\begin{tabularx}{\columnwidth}{
    >{\centering\arraybackslash}p{0.24\columnwidth}
    >{\raggedright\arraybackslash}X
}
\hline
\hl{\textbf{Notation}} & \hl{\textbf{Description}} \\
\hline
\hl{$R$} &
\hl{Number of hash functions used to generate blocking keys for each entity.} \\

\hl{$\theta$} &
\hl{Collision-frequency threshold used to promote a candidate to a matching pair.} \\

\hl{$\Delta$} &
\hl{Temporal eviction window, expressed in discrete processing increments.} \\

\hl{$t_i$} &
\hl{Discrete timestamp assigned when entity $e_i$ is inserted into the blocking structure; smaller values indicate older entities.} \\

\hl{offsetA} &
\hl{Number of Dataset 1 entities processed in each Blocking batch.} \\

\hl{offsetB} &
\hl{Number of Dataset 2 entities processed in each Matching batch.} \\

\hl{Candidate pair} &
\hl{Pair of entities retrieved because they co-occur in one or more approximate blocks.} \\

\hl{Matching pair} &
\hl{Candidate pair whose collision frequency satisfies the threshold $\theta$.} \\

\hline
\end{tabularx}
\end{table}


\section{Related Work} \label{sec:trabRelacionados}

In the literature, several ER solutions have been proposed over time, each adopting different techniques to address diverse application scenarios. Online techniques, such as those described in \cite{DBLP:conf/edbt/KarapiperisGV18, 10.14778/2556549.2556567}, aim to handle real-time data through incremental updates and are useful in applications requiring low-latency updates, such as e-commerce and web platforms. Progressive techniques, such as those presented in \cite{DBLP:conf/edbt/GazzarriH23, 8403302}, seek to optimize ER by prioritizing the identification of the most promising matches, providing faster partial results even when dealing with large volumes of heterogeneous data. Incremental techniques \cite{10.14778/2732939.2732943, DONASCIMENTO2018335, inproceedingssuperblocking} are designed to handle the continuous arrival of data by efficiently updating matching results as new information is integrated into the system and previously committed errors are corrected. Topic-aware techniques \cite{10.1145/3448016.3457238} leverage data context to improve the entity matching task. Finally, several works propose ER techniques for streaming environments \cite{info13120568, 9378127, Karapiperis2018SummarizationAF}, targeting the processing of continuously arriving, high-volume data while addressing challenges related to memory management and real-time processing efficiency.

Research focused on ER in streaming data environments proposes techniques to improve performance under time and memory constraints. Works such as ExpBlock \cite{karapiperis2023} employ a block eviction strategy based on the random selection of candidate blocks in order to free memory when predefined block storage limits are reached. In addition, the authors propose an entity cleaning process within blocks using the Bernoulli process \cite{IBE2014369}, a probabilistic model based on independent experiments with two possible outcomes: success or failure. In this context, the process is applied through binary variables, where each entity in a block is associated with a probability of being removed (success) or retained (failure), enabling the elimination of irrelevant entities within blocks. Experimental results show that ExpBlock, in most cases, achieves performance comparable to or better than other approaches, such as \cite{Gazzarri2021} and \cite{10.1145/3341105.3375776}, when considering Recall and Precision metrics.

A limitation of ExpBlock \cite{karapiperis2023} is that related entities arriving with large time gaps may be incorrectly grouped, making the identification of matches more difficult. In such scenarios, the Recall metric may be negatively affected, as blocks may be evicted before all related entities have arrived, prioritizing memory release and reduced processing time at the expense of identifying all relevant entity correspondences.

In another work, the authors of \cite{10.1109/TKDE.2024.3361022} propose the AB blocking technique, which is based on the use of hash-based representations derived from strings to generate entity block keys. This technique relies on theoretical guarantees to group entities, as demonstrated in \cite{lshbased01}. The entity blocking process involves generating a fixed number of distinct hash codes based on entity attribute values and checking intersections between the generated values and those of other entities to determine matches. This approach simplifies the ER blocking stage, as it avoids complex operations for entity pair detection.

In \cite{info13120568}, the authors propose a blocking technique capable of handling incremental streaming data in a \textit{schema-agnostic} manner. The technique selects entity attributes while evicting those that may negatively affect the blocking stage and employs algorithms tolerant to data noise. These algorithms generate hash values from attribute values, supporting the effective creation of entity blocks in which entities share similar or identical attributes. In addition, the authors introduce an incremental and parallel pipeline for entity block generation, designed to efficiently handle streaming data.

The work in \cite{karapiperis2023} leverages both block cleaning and random selection strategies to improve the efficiency of the blocking stage, whereas the approach proposed in \cite{info13120568} stands out by combining attribute selection and entity block cleaning to address streaming data, maintaining matching quality under memory constraints and continuous data arrival. In contrast, the proposal in \cite{10.1109/TKDE.2024.3361022} primarily focuses on random entity selection to generate correspondences at low computational cost and higher speed, without requiring block cleaning procedures or attribute selection.

However, these solutions \cite{10.1109/TKDE.2024.3361022, karapiperis2023, info13120568, lshbased01, Gazzarri2021, 10.1145/3341105.3375776, DBLP:conf/edbt/KarapiperisGV18, Karapiperis2018SummarizationAF, 10.14778/2556549.2556567, DBLP:conf/edbt/GazzarriH23, 8403302, DONASCIMENTO2018335, 10.1145/3448016.3457238, 9378127, 10.14778/2732939.2732943} address ER in a traditional manner, where block construction and access to data structures are treated as consolidated stages in the literature.

Recent works have started to broaden this perspective by considering fairness and explainability in ER \cite{11129015}, as well as by proposing auto-configurable ER pipelines that adapt to different data scenarios \cite{11129641}. The first work provides a systematic overview of how fairness constraints and explanation techniques can be integrated into the different stages of the ER pipeline, while the second proposes an AutoML-based framework that automatically selects language models, blocking parameters, and clustering strategies to optimize ER performance across diverse datasets. Although these approaches do not directly target eviction strategies for in-memory blocks, they highlight the need for more holistic ER designs that consider not only performance, but also transparency, adaptability, and user trust.

On the other hand, Deep Learning–based approaches \cite{kasai2019lrder, dev2022DeepLI, Li_2020, zeakis2023pretrainedembeddingsentityresolution} incorporate Deep Learning into the ER process, differing in terms of the adopted architectures, processing pipelines, and the level of integration between blocking and matching.

\textit{Ditto} \cite{Li_2020} reformulates Entity Matching as a text-pair classification task using \textit{Transformer-based} models such as BERT and RoBERTa. The technique employs fine-tuning, knowledge injection, and data augmentation, achieving expressive gains of up to 29\% in \textit{F1-score} even with limited labeled data. However, its pairwise processing strategy and high computational cost limit its scalability in large-scale or real-time scenarios.

In \cite{zeakis2023pretrainedembeddingsentityresolution}, a comparative analysis of 12 pre-trained language models is conducted across the vectorization, blocking, and \textit{matching} stages, considering both \textit{schema-aware} and \textit{schema-agnostic} settings. The results reveal significant trade-offs among computational cost, robustness, and effectiveness, highlighting scalability challenges and guiding model selection according to the application domain and the specific stage of the ER pipeline.

The work of \cite{dev2022DeepLI} proposes a \textit{DistilBERT}-based approach combined with active learning for streaming ER. The technique supports incremental training and classification in streaming environments, achieving \textit{F1-scores} of up to 0.97. Despite its strong performance, the reliance on continuous labeling and the associated computational cost remain important challenges.

Finally, \cite{10.1145/3786768} introduces the CRL (\textit{Consistent Record Linkage}) metric, which enables the evaluation of ER classifiers based on user-defined error thresholds. The approach demonstrates strong performance in Big Data scenarios by balancing precision, \textit{recall}, and computational efficiency without requiring fully labeled datasets.

\hl{In contrast, the present work proposes eviction strategies for managing in-memory blocks in streaming ER and evaluates their measured effects on processing cost and matching quality without assuming that quality is preserved. The evaluated criteria are entity age and block capacity.} Table~\ref{tab:trabalhos} presents a comparative overview of existing techniques across four dimensions: Heterogeneity, Noisy Data, Streaming, and Incremental Processing. Each row corresponds to a group of related works, while each column indicates whether a given technique addresses a specific research aspect. An “X” denotes support for the corresponding characteristic. The table highlights that only a few approaches comprehensively cover all dimensions, revealing a gap in the literature that motivates the contributions of the proposed work.

Among existing approaches, this work focuses on traditional ER techniques in \textit{streaming} environments, which follow the classical ER pipeline based on blocking and explicit attribute comparisons. These techniques are favored due to their lower computational requirements and more predictable performance. Moreover, they are well established in the literature, featuring mature and extensively evaluated techniques with structured blocking and indexing mechanisms for entities and candidate blocks, which enhances reproducibility and experimental extensibility. By adopting a traditional approach, this study seeks to balance efficiency and accuracy, ensuring the scalability required for large-scale continuous data processing while maintaining low computational cost and reduced architectural complexity.

\begin{table}[!t]
\centering
\caption{Overview of related work. "X" indicates that the work addresses the research topic whilst blank cell indicates that it does not.}
\label{tab:trabalhos}
\scriptsize
\setlength{\tabcolsep}{2pt}
\renewcommand{\arraystretch}{0.95}
\begin{tabularx}{\columnwidth}{>{\raggedright\arraybackslash}p{0.36\columnwidth} >{\centering\arraybackslash}p{0.14\columnwidth} >{\centering\arraybackslash}p{0.14\columnwidth} >{\centering\arraybackslash}p{0.14\columnwidth} >{\centering\arraybackslash}p{0.14\columnwidth}}
\hline
\textbf{Articles} & \textbf{Heterog.} & \textbf{Noisy} & \textbf{Streaming} & \textbf{Increm.} \\
\hline \cite{DBLP:conf/edbt/KarapiperisGV18, karapiperis2023, Karapiperis2018SummarizationAF} &  &  & X &  \\
\hline \cite{DONASCIMENTO2018335, 10.14778/2732939.2732943,inproceedingssuperblocking} &  &  &  & X \\
\hline \cite{DBLP:conf/edbt/GazzarriH23, 8403302} & X &  &  & X \\
\hline \cite{10.1145/3448016.3457238, dev2022DeepLI} &  &  & X & X \\
\hline \cite{10.14778/2856318.2856326, 10.1145/3297280.3299730, Papadakis2016ScalingER} & X &  &  &  \\
\hline \cite{zeakis2023pretrainedembeddingsentityresolution, lshbased01} &  & X &  &  \\
\hline \cite{Li_2020} & X & X &  &  \\
\hline \cite{info13120568} & X & X & X & X \\
\hline \cite{kasai2019lrder} &  &  &  &  \\
\hline \cite{9378127} & X & X & X &  \\
\hline \cite{10.1109/TKDE.2024.3361022} & X &  & X & X \\
\hline Proposed Work & X & X & X & X \\
\hline
\end{tabularx}
\end{table}

\section{Entity and Block Eviction Strategies} \label{sec:estrategias-descarte-entidades}

In this section, we describe the operation of the different Entity Eviction techniques proposed in this work: (i) Entity Eviction technique; (ii) Entity Block Eviction technique; and (iii) Global Entity Eviction technique. This set of techniques was designed in a complementary manner, since each eviction technique relies on different criteria, allowing them to be combined and to address distinct aspects of block management throughout the ER task. However, in order to isolate the individual impact of each technique and enable a more controlled and interpretable experimental analysis, this work evaluates each eviction technique independently, without exploring their combined application. This methodological choice allows for a clearer understanding of the specific effects of each Eviction mechanism on both memory usage and the quality of the generated matches.

In the first technique, the lifetime of entities within a block is used as the Eviction criterion. The underlying assumption is that older entities tend to contribute less to the ER task, as they may become obsolete or low-discriminative over time.

In the second technique, Eviction occurs at the block level, meaning that when a stored entity block reaches a predefined capacity limit (expressed as the maximum number of entities allowed in that block) the entire block is removed from memory. This approach is based on the hypothesis that excessively large blocks are typically formed from very generic block keys, such as “gender” or “country of origin,” which may not adequately describe the entities. As a result, such blocks tend to exhibit low quality and high memory consumption.

\hl{Finally, the third evaluated technique is a hybrid Global Eviction policy: it combines age-based removal across an insertion-order queue with local block clearing at a capacity trigger. The implementation does not use access frequency or matching contribution as eviction criteria.} Next, we detail the operation of each proposed Eviction technique.

\subsection{Entity Eviction} \label{subsec:descarte-entidades}

A removal policy based solely on individual entity eviction, without considering the time at which entities are inserted into blocks, ultimately ignores the temporal context of the entities, as shown in \cite{10.1109/TKDE.2024.3361022}. In this sense, this research proposes an improvement to the AB entity block structure with the goal of considering the dynamics of new entity arrivals and their impact on the ER task over time. Initially, it is necessary to incorporate the instant at which an entity is processed during the blocking stage, denoted as $t_i$. With the introduction of this temporal marker, it becomes possible to track the exact time an entity enters a block, thereby enabling the removal of entities that arrived within a given time interval.

Formally, let \( \mathcal{E} = \{e_1, e_2, \dots, e_n\} \) be the set of processed entities, where each entity \( e_i \) is represented by an ordered pair \( e_i = (id_i, t_i) \), in which \( id_i \) corresponds to the unique identifier of the entity and \( t_i \in \mathbb{N}^{+} \) denotes the discrete processing time at which the entity was inserted into the blocking structure. Thus, entities with smaller values of \( t_i \) correspond to older records in the stream, whereas larger values indicate more recently processed entities.

Let \( \mathcal{B} = \{B_1, B_2, \dots, B_m\} \) denote the set of entity blocks, where each block \( B_j \subseteq \mathcal{E} \) stores entities that share the same blocking key. Let \( L \) be the maximum number of entities allowed in each block, that is,
\[
|B_j| \leq L, \quad \forall B_j \in \mathcal{B}.
\]

Let \( \Delta \in \mathbb{N}^{+} \) denote the temporal eviction window, which defines the maximum allowed lifetime of entities within a block. Based on the current processing time \( t \), the eviction threshold is defined as
\[
t_d = t - \Delta,
\]
where entities with timestamps smaller than or equal to \( t_d \) are considered obsolete and therefore eligible for removal.

When a block \( B_j \) reaches the maximum limit \( L \), an eviction process is triggered based on the temporal threshold \( t_d \), removing all entities in the block whose timestamps satisfy this condition. The set of entities removed from block \( B_j \) is defined as
\[
B_{j}^{r} =
\{ e_i \in B_j \mid t_i \leq t_d \},
\quad \text{with } |B_j| > L.
\]

This technique ensures that each block maintains a controlled size, freeing space for the insertion of new entities and preventing unbounded growth of the blocking structure. It introduces a temporal prioritization policy, whereby older records tend to be removed in favor of more recent entities that are potentially more relevant to the ER task.

Unlike the approach presented in \cite{10.1109/TKDE.2024.3361022}, which evict exclusively the oldest entity in each block, this work adopts a policy based on a temporal window, allowing the simultaneous removal of multiple obsolete entities. This strategy improves the efficiency of the eviction process and enables more robust handling of bursts of entity insertions that occur within short time intervals.

Analogous to cache mechanisms in web systems, this policy automatically removes old or low-relevance data, ensuring that the most recent information with greater potential to contribute to the identification of correspondences remains available in memory.

\begin{figure}[h]
  \centering
  \includegraphics[width=0.9\linewidth]{images/descarte_entidades_1.png}
  \caption{Example of entity eviction at time instant $t_5$.}
  \label{fig:descarte_entidades}
\end{figure}

Figure~\ref{fig:descarte_entidades} illustrates an example of the entity eviction technique. Entity $e_9$ is processed at time instant $t_5$ and must be allocated to the blocks with blocking keys $-5634$, $-6583$, $-7645$, and $-2452$, i.e., blocks $b_7$, $b_9$ (created at the time of insertion), $b_6$, and $b_2$, respectively. However, due to the block capacity limit (defined in this example as four entities) the technique removes all entities whose arrival time is earlier than or equal to $t_1$ in order to free memory space. It can be observed that entities $e_1$ and $d_3$, inserted at time instant $t_1$, are removed. It should be emphasized that the technique operates only on blocks that have already reached their maximum capacity. Algorithm \ref{alg:blocking_array_simplified} describes the operation of the entity eviction technique.

\begin{algorithm}[H]
\caption{Time-aware entity eviction}
\label{alg:blocking_array_simplified}
\begin{algorithmic}[1]

\Require Set of entities $\mathcal{E}$; set of hash functions $\mathcal{H}$; block capacity limit $L$; eviction window size $\Delta$
\Ensure Block structure $\mathcal{B} = {\mathcal{B}_1, \dots, \mathcal{B}_L}$

\State Initialize block structure $\mathcal{B}$ and metadata structure $\mathcal{S}$ \Comment{Stores timestamps per block}
\State $t \gets 0$ \Comment{Global time counter}

\For{each $e_i \in \mathcal{E}$}
\For{each $h_l \in \mathcal{H}$}
\State Compute block key $k \gets h_l(e_i)$
 \If{block $b_{l,k}$ exists}
        \If{$|b_{l,k}| \geq L$}
            \State $t_d \gets t - \Delta$ \Comment{Temporal eviction threshold}
            \State Remove all entities $e_j \in b_{l,k}$ such that $t_j \leq t_d$
            \Comment{Eviction outdated entities}
        \EndIf
        \State $b_{l,k} \gets b_{l,k} \cup \{(e_i,t)\}$
    \Else
        \State $\mathcal{B} \gets \mathcal{B} \cup \{b_{l,k}\}$
        \State $b_{l,k} \gets b_{l,k} \cup \{(e_i,t)\}$
    \EndIf
\EndFor
\State $t \gets t + 1$
\EndFor
\end{algorithmic}
\end{algorithm}

\subsection{Block Eviction} \label{subsec:descarte-blocos}

When blocks frequently reach their capacity limit, it is likely that the entities contained in these blocks include low-discriminative data, such as \textit{stopwords}\footnote{Stopwords are common words with little semantic value, typically removed in NLP tasks, such as ``the'', ``a'', ``of'', and ``in''.}, which do not significantly contribute to determining similarity between entity pairs.

Given this limitation, an alternative strategy is to evict entire blocks once they reach a predefined maximum capacity. Although this threshold is selected in this work based on practical considerations, its value can be adjusted according to the characteristics of the ER task and the underlying data.

The proposed technique for full block eviction aims to optimize memory usage and processing efficiency. Block eviction is triggered when a block reaches its capacity limit. Let \( b_j = \{ e_1, e_2, \dots, e_k \} \) be a block containing \( k \) entities, where \( L \) denotes the fixed maximum number of entities allowed per block. To ensure efficient memory management and adequate performance of the ER task, it must be guaranteed that \( k \leq L \) for all blocks \( b_j \) in the block set \[B = \{ b_1, b_2, \dots, b_n \}.\]

It is important to note that the threshold \(L\) adopted in this strategy corresponds to the same block capacity parameter employed in the Entity Eviction approach (Section~\ref{subsec:descarte-entidades}). However, its role differs between the two techniques. In Entity Eviction, \(L\) acts as a trigger for temporal maintenance operations, causing the selective removal of obsolete entities while preserving the block itself. In contrast, in Block Eviction, \(L\) defines a hard capacity limit whose violation results in the removal of the entire block.

When \( k = L \), a block \( b_j \) is evicted from \( B \), formally expressed as:
\[
\exists b_j \in B \, , \ (|b_j| = L) \Rightarrow B' = B - \{ b_j \}.
\]

The set of remaining blocks after the eviction of \( b_j \) is denoted by \( B' \), and it must satisfy the following condition:
\[
\forall b_j \in B', \ |b_j| \leq L.
\]
In this way, it is ensured that the block set \( B \) contains only blocks whose number of entities does not exceed the limit \( L \), thereby optimizing resource usage.

\begin{figure}[h]
  \centering
  \includegraphics[width=0.9\linewidth~]{images/descarte_blocos_1.png}
  \caption{Example of block eviction.}
  \label{fig:descarte_blocos}
\end{figure}

By evicting oversized blocks, the system reduces computational overhead and accelerates entity processing. In systems that handle data streams, as in the context of this work, this approach is even more crucial. The block eviction process therefore not only improves storage efficiency but also frees up space for the insertion of more relevant data.

This process ensures that only pertinent information is preserved, resulting in increased precision and relevance of the comparisons. Through the controlled eviction of unnecessary entities, the task preserves essential data and avoids the removal of entities that could compromise the accuracy of the ER task.

Figure~\ref{fig:descarte_blocos} illustrates the operation of the Block Eviction technique. At a given instant in the data stream processing (\(t_9\)), a new entity, represented by \(e_9\), is received and must be allocated to block \(b_2\). However, since the block has already reached its maximum storage capacity (four entities), the block is entirely evicted to enable the insertion of the new entity. In this process, the Block Eviction technique removes all entities previously stored in block \(b_2\) and inserts the new entity (\(e_9\)) into \(b_2\). Algorithm \ref{alg:blocking_general} presents the operation of the Block Eviction technique.

\begin{algorithm}[H]
\caption{Block Eviction Strategy}
\label{alg:blocking_general}
\begin{algorithmic}[1]
\Require Set of entities $\mathcal{E}$; set of hash functions $\mathcal{H}$; block capacity limit $L$; eviction window size $\Delta$
\Ensure Block structure $\mathcal{B} = {\mathcal{B}_1, \dots, \mathcal{B}_L}$
\State Initialize block structure $\mathcal{B}$ and metadata structure $\mathcal{S}$ \Comment{Stores timestamps per block}
\State $t \gets 0$ \Comment{Global time counter}
\For{each $e_i \in \mathcal{E}$}    
    \For{each $h_l \in \mathcal{H}$}
        \State Compute the blocking key $k \gets h_l(e_i)$   
        \If{block $b_{l,k}$ exists}
             \If{$|b_{l,k}| \geq L$}
                \State Remove the entire block $b_{l,k}$ \Comment{Whole-block eviction when capacity is reached}
                \State Create a new empty block $b_{l,k}$        
            \EndIf
            \State $b_{l,k} \gets (e_i, t)$
        \Else
            \State $\mathcal{B} \gets b_{l,k}$ \Comment{Create new block}
            \State $b_{l,k} \gets b_{l,k} \cup \{(e_i, t_i)\}$
        \EndIf
    \EndFor
    \State $t \gets t + 1$
\EndFor
\end{algorithmic}
\end{algorithm}

\subsection{Global Entity Eviction} \label{subsec:descarte-entidades-global}

The third technique proposed in this work consists of a global entity eviction technique. \hl{The evaluated implementation is a hybrid policy: before insertion, it clears an existing block when that block contains 1,000 entities, and it also globally expires entities outside a 2,500-entity retention window. Thus, its global temporal rule is combined with a local capacity rule.}

\hl{For each inserted Dataset 1 entity, the implementation records its identifier, insertion increment, and associated block keys in a FIFO queue. Once the insertion increment exceeds the retention window, queue entries older than the current cutoff are removed, along with their identifiers from the associated blocks when still present at the front of those blocks.}

\hl{The temporal criterion is an age-based heuristic and does not determine whether a record is semantically irrelevant or has already contributed to a match. A record outside the configured window may still be relevant to a later match.}

\hl{The experiment evaluates global expiration and local block clearing together; it does not include an ablation that isolates either component. Consequently, the separate contribution of each rule cannot be inferred from these results.}

\hl{The global rule uses each entity's insertion increment as its expiration criterion. It does not estimate the entity's probability of contributing to the Matching stage.}

\hl{The FIFO queue $\mathcal{Q}$ stores one entry $(id_i,t_i,K_i)$ per inserted Dataset 1 entity, where $t_i$ is its insertion increment and $K_i$ is the set of block keys receiving that entity. Entries are appended in stream order and examined from the oldest end when the retention cutoff advances.}

\begin{figure}[h]
  \centering
  \includegraphics[width=0.9\linewidth~]{images/descarte_global.png}
  \caption{Example of global entity eviction.}
  \label{fig:fluxo_heap}
\end{figure}

\hl{Figure~\ref{fig:fluxo_heap} illustrates the policy with a toy retention window of 10 increments: when $e_{11}$ arrives, entities outside that window are removed from their associated blocks. The evaluated ``descarte velho'' run uses a 2,500-entity window, not the illustrative value of 10.}

\hl{Algorithm~\ref{alg:blocking_heap} formalizes the evaluated hybrid policy. Global expiration and local block clearing are evaluated together; no ablation is available to attribute their effects separately.}

\begin{algorithm}[H]
\caption{Entity Eviction with Global Control}
\label{alg:blocking_heap}
\begin{algorithmic}[1]
\Require Set of entities $\mathcal{E}$; set of hash functions $\mathcal{H}$; \hl{capacity trigger $L=1000$; retention window $\Delta=2500$}
\Ensure \hl{Block structure $\mathcal{B}$}
\State \hl{Initialize empty block structure $\mathcal{B}$ and FIFO queue $\mathcal{Q}$; set $t\gets0$}
\For{each $e_i \in \mathcal{E}$}
    \State \hl{$K_i\gets\emptyset$}
    \For{each $h_l \in \mathcal{H}$}
        \State Compute the blocking key $k \gets h_l(e_i)$
        \If{block $b_{l,k}$ exists}
            \If{$|b_{l,k}| \geq L$}
                \State \hl{Clear $b_{l,k}$ (local capacity clearing)}
            \EndIf
        \Else
            \State \hl{Create an empty block $b_{l,k}$}
        \EndIf
        \State \hl{Insert $e_i$ into $b_{l,k}$}
        \State \hl{$K_i\gets K_i\cup\{(k,l)\}$}
    \EndFor
    \State \hl{Append $(id_i,t,K_i)$ to $\mathcal{Q}$}
    \If{$t>\Delta$}
        \State \hl{Set $t_{min}\gets t-\Delta$}
        \While{$\mathcal{Q}$ is not empty and the oldest entry has timestamp $<t_{min}$}
            \State \hl{Remove the oldest entry $(id_j,t_j,K_j)$ from $\mathcal{Q}$}
            \For{each $(k,l)\in K_j$}
                \State \hl{Remove $id_j$ from the front of $b_{l,k}$ if it is still there}
            \EndFor
        \EndWhile
    \EndIf
    \State \hl{$t\gets t+1$}
\EndFor
\end{algorithmic}
\end{algorithm}

\section{Experiments} \label{sec:avaliacao}

This section presents the experimental evaluation and discusses the results. The experiments were mainly intended to analyze the behavior of the ER task using the proposed eviction techniques. As the baseline for all experiments, the AB blocking technique \cite{10.1109/TKDE.2024.3361022}, discussed in Section \ref{sec:trabRelacionados}, was used. It is important to highlight that the AB technique is limited to entity eviction without considering the time at which an entity is inserted into a block, relying only on the removal of the first entity inserted into the block. The evaluation focused on the Recall and Precision metrics. The values of these metrics range in the interval [0, 1], where higher values indicate better performance. In addition, the execution time of the Blocking and Matching stages was also evaluated.

Although the primary focus of this work is the blocking stage, we also account for the matching stage in AB, which compares entities within each block to identify matching pairs. Thus, the Recall and Precision metrics are computed based on the pairs identified as matches in this stage.

\hl{In the baseline and Experiments 2 and 3 (Block Eviction), each Dataset 2 record is processed in row order using the same selected attributes and $q$-gram MinHash procedure used for blocking. For each record, the matcher draws $R'$ hash-function indices independently and uniformly from the $R$ available indices. For each sampled index, it retrieves identifiers stored under the corresponding block key and increments a collision counter. An identifier becomes a candidate when its counter reaches $\lceil \theta R' \rceil$; with $R'=50$ and $\theta=0.5$, this requires 25 collisions. The emitted pairs are compared with the ground-truth mapping to classify $TP$ and $FP$; no additional attribute-level similarity function is applied. The ``descarte velho'' run differs: it uses ground truth to select Dataset 2 query records, as detailed in Experiment 3, so its quality outputs are not directly comparable.}

\hl{Let $TP$ denote emitted candidate pairs present in the ground truth, $FP$ emitted pairs absent from the ground truth, and $FN$ ground-truth pairs not emitted. Recall and Precision are $TP/(TP+FN)$ and $TP/(TP+FP)$, respectively; for the baseline and Experiments 2 and 3, $FN=|MP|-TP$ and the candidate-pair count equals $TP+FP$. The full Cartesian product is not enumerated, so $TN$ is not measured and a meaningful true-negative count is unavailable. Table~\ref{tab:resultados_consolidados_rotacionada} reports $TP$, $FP$, and $FN$ by dataset and method. The raw ``descarte velho'' counters do not satisfy this accounting and are presented separately as diagnostics, not as a confusion matrix.}

An additional experiment was conducted to evaluate the impact of the volume of data processed simultaneously in the ER task. For this purpose, the \textit{offset} parameter was used, which indicates the size of the data blocks analyzed at each stage. This parameter was applied to both the Blocking and Matching stages and varied according to a geometric progression, in order to cover different processing scales. This variation made it possible to measure how different volumes of processed entities affect the performance of the evaluated technique, both in execution time and in the quality of the identified entity pairs, providing evidence to assess its scalability. The choice to employ a single technique aimed to isolate the impact of the \textit{offset} parameter, avoiding interference arising from structural or heuristic differences among multiple approaches.

\hl{Experiments 2 and 3 use a block-capacity trigger of 1,000 entities and process 50 records per batch in Blocking and Matching. In Experiment 2, each block starts with an insertion-time cutoff of 50; whenever that block reaches capacity, entries older than its cutoff are removed from its front and the cutoff advances by 50. This cutoff advances on capacity-triggered checks, even if no entries are removed, so 1,000 is not a guaranteed hard bound for this policy. Experiment 2's Block Eviction clears the whole block at the 1,000-entity trigger. The Global Eviction run named ``descarte velho'' combines local clearing at 1,000 entities with a global retention window of 2,500 processed entities. The 1,000-entity trigger follows the initial capacity used by the cited AB configuration; the 50-entity cutoff increment matches the batch size. The 2,500-entity global window is a fixed implementation setting, not a tuned or validated optimum.}

\subsection{Environment and Configuration}
\hl{The declared runtime dependencies are mmh3 4.0.1, NumPy 1.24.4, pandas 2.0.3, and SciPy (version 1.9.0 or later), as listed in the project requirements file.}
The experimental environment consisted of a physical machine with 32 GB of RAM operating at 3{,}000 MHz, and an Intel Core i5-10400 processor @ 2.90 GHz with 6 cores. The operating system used was Ubuntu 22.04.4 LTS, 64-bit. The development environment was configured with Conda version 23.11.0 and Python 3.11.5. \hl{The V2 comparisons reported here use 10 replicas with consecutive seeds from 20260907 through 20260916; reported summary values are replica means.}\footnote{The link to the project repository will be made available after the blind-review stage.}

\hl{For the baseline and Experiments 2 and 3, the stream order is the row order of Dataset 1 for blocking and Dataset 2 for matching, processed in batches of 50 without shuffling. Each replica reinitializes the pseudo-random generator with its stated seed. The implementation uses $q=2$, $R=116$, $R'=50$, and $\theta=0.5$. The ``descarte velho'' run is an exception: it does not process Dataset 2 independently in row order, because its queries are selected through the ground-truth mapping, as described in Experiment 3.}

\hl{For each configuration, the scripts persist replica-level values for blocking time, matching time, $TP$, $FP$, candidate-pair count, Recall, and Precision. They report the sample mean, sample standard deviation, and a two-sided 95\% confidence interval based on the Student $t$ distribution. In the result tables, each cell is shown as the replica mean $\bar{x}$ plus or minus the confidence-interval half-width $h=t_{0.975,n-1}s/\sqrt{n}$; the corresponding interval is $[\bar{x}-h,\bar{x}+h]$. For these runs, $n=10$ and the interval uses 9 degrees of freedom. Values are rounded for display, with intervals calculated from unrounded replica values. The V2 comparisons of Experiments 2 and 3 with the baseline use two-sided Welch independent-samples $t$-tests with unequal variances; although the replicas use the same seed values, the scripts do not perform paired tests. The tests are exploratory and are not adjusted for multiple comparisons. Constant metric samples produce zero-width intervals and undefined test statistics, reported as n/a in the test table. The ``descarte velho'' Global Eviction quality metrics are excluded from these tests because its matching input is selected using ground truth, as detailed in Experiment 3.}

\subsection{Threats to Validity}
\hl{Internal validity is limited by the fixed stream order, the synthetic perturbation procedure used to create the second NCVR and MusicBrainz sources, and the randomized sampling of hash indices. Fixed seeds and repeated runs reduce run-to-run variation but do not eliminate dependence on the selected ordering or perturbation settings. Construct validity depends on the correctness of the ground-truth mappings and the implementation of $TP$, $FP$, and $FN$; therefore, all quality-related interpretations remain conditional on validation of these calculations.}

\hl{External validity is limited to the three evaluated datasets, the selected attributes, the configured parameter values, and execution on a single-machine environment. The study does not evaluate distributed deployment, online concept drift, or production arrival processes; consequently, the results should not be generalized to those settings without additional experiments.}

\subsection{Datasets and Characteristics}
Three distinct datasets were used in the experiments: Scholar-DBLP (SCD)\footnote{https://dbs.uni-leipzig.de/research/projects/benchmark-datasets-for-entity-resolution}, which contains bibliographic records in the Computer Science domain; NCVR\footnote{https://www.ncsbe.gov/results-data/voter-registration-data}, which comprises voter registration records from the U.S. state of North Carolina; and MusicBrainz\footnote{https://dbs.uni-leipzig.de/research/projects/benchmark-datasets-for-entity-resolution}, which contains records derived from the MusicBrainz dataset\footnote{https://musicbrainz.org/}.

The SCD dataset consists of two distinct files, denoted as A and B, representing two different data sources. In this case, file A is used in the blocking stage, while file B is employed in the entity matching stage.

\hl{The ground-truth pairs are loaded from 	exttt{truth.csv} (ID columns 	exttt{idDBLP}/	exttt{idScholar}), 	exttt{ground	extunderscore truthNC	extunderscore voter.csv} (	exttt{antigoId}/	exttt{novoId}), and 	exttt{ground	extunderscore truth	extunderscore music	extunderscore brainz.csv} (	exttt{antigoidmusic2}/	exttt{novoidmusic1}), respectively. The evaluation denominators $|MP|$ are 5,347 for SCD, 40,893 for NCVR, and 193,471 for MusicBrainz; the scripts set the first two values explicitly and use the MusicBrainz ground-truth file row count.}

In contrast, the NCVR and MusicBrainz datasets consist of a single file, characterizing a single-source scenario. To enable the matching stage for these datasets, a second file (B) was artificially generated from the original dataset. This process was carried out using a Python script, which allowed splitting the records into two distinct files while maintaining consistency with the adopted experimental protocol. The second file was constructed according to the rules described in \cite{karapiperis2023}, in which each entity has a probability of being perturbed, generating duplicate entities with slight variations in their attribute values. These variations consisted of, for NCVR, adding 1 year to the birth year of the selected entity and converting the entity's last name to lowercase. For MusicBrainz, an analogous procedure was applied: 1 was added to the song's year, while the artist's name was fully converted to lowercase.

The three datasets present different structures and attributes that are fundamental for the ER task. The SCD dataset mainly contains the properties \textit{id}, \textit{title}, \textit{authors}, \textit{venue}, and \textit{year}, which describe identifiers, publication titles, authors, venue, and year of publication. The NCVoter dataset includes a broader set of attributes, among which \textit{ncid}, \textit{first\_name}, \textit{last\_name}, \textit{registr\_dt}, and \textit{age\_at\_year\_end} stand out. Although the dataset contains additional fields, these represent the essential information related to voter identity and characteristics. Similarly, the MusicBrainz dataset includes a large number of properties, but the most relevant for this study are \textit{title}, \textit{artist}, \textit{album}, \textit{year}, and \textit{language}, which capture core information about songs and their metadata. These selected attributes reflect the main variables required to conduct consistent blocking and entity matching experiments.

Table \ref{tab:tab1} presents the characteristics of the datasets used. The rows $\lvert A \rvert$ and $\lvert B \rvert$ correspond, respectively, to the number of entities in the datasets used for blocking and those used for matching. The row $\lvert MP \rvert$ represents the number of true matching pairs between entities in the datasets, while $\lvert AT \rvert$ indicates the number of available attributes. For instance, in the SCD dataset, the first file (A) contains 64 thousand entities, while the second file (B) contains 2,600 entities. The number of matching entity pairs between the two datasets is 5 thousand. Both datasets A and B contain 5 attributes.

\begin{table}[ht]
\centering
\caption{Characteristics of the datasets.}
\begin{tabular}{|c|c|c|c|}
\hline
& \textbf{SCD} & \textbf{NCVR} & \textbf{MusicBrainz} \\ \hline
    \textbf{$|A|$} & 64k  & 40k  & 190k \\ \hline
    \textbf{$|B|$} & 2.6k & 40k  & 190k \\ \hline
    \textbf{$|MP|$} & \hl{5,347}   & \hl{40,893} & \hl{193,471} \\ \hline
    \textbf{$|AT|$ }  & 5    & 58 & 12   \\ \hline
\end{tabular}
\label{tab:tab1}
\end{table}


The transformation of entity attribute values into hash representations was performed using the MinHash algorithm\footnote{https://github.com/jmhodges/minhash}. Each entity was assigned to blocks based on keys derived from the generated hash values. In the MinHash process, multiple hash functions are applied to the elements composing the entity's attribute set, producing integer values; for each function, the minimum value obtained is selected as a component of the MinHash signature. The attributes used for this algorithm were as follows: \texttt{title} and \texttt{authors} for the SCD dataset; \texttt{first\_name}, \texttt{last\_name}, \texttt{registr\_dt}, and \texttt{age\_at\_year\_end} for the NCVR dataset; and \texttt{title}, \texttt{length}, \texttt{artist}, and \texttt{album} for the MusicBrainz dataset.

\subsection{Experiment 1: Varying the Number of Entity Evictions}
In order to analyze the influence of the number of entities removed during the eviction process, an experiment was conducted using the SCD dataset, in which the number of evicted entities increased progressively at each cycle, as illustrated in Figure \ref{fig:recall_100-500}.

\begin{figure}[h]
  \centering
  \includegraphics[width=0.9\linewidth]{images/va.png}
  \caption{Evolution of Recall and Precision during the eviction cycles.}
  \label{fig:recall_100-500}
\end{figure}

As can be observed, there was no significant improvement in the Recall and Precision metrics when the number of entities removed in each eviction process was increased. This can be explained by the fact that blocks reaching their maximum capacity tend to store information that does not contribute to the quality of the results, potentially containing data that are irrelevant to the ER task. The results obtained in this experiment provided a solid basis for proposing a more efficient eviction process, which does not rely solely on the removal of individual entities, but rather on the elimination of entire blocks of entities.

\subsection{Experiment 2: Evaluating the Entity Eviction and Block Eviction Techniques}
\hl{This section compares the AB baseline with Entity Eviction and Block Eviction on the three datasets using the V2 replica outputs. Blocking time, matching time, candidate-pair count, and the confusion-matrix components are summarized in Table~\ref{tab:resultados_consolidados_rotacionada}. The reported candidate-pair count satisfies $TP+FP$ in these runs. False negatives are calculated as the number of ground-truth pairs minus $TP$; true negatives are not enumerated because the full Cartesian product of each dataset pair is not evaluated.}

\hl{On SCD, Entity Eviction has a higher mean matching time than the baseline (216.571~s versus 155.856~s) and nearly unchanged candidate-pair count (74,596.4 versus 75,527.1). Block Eviction reduces matching time to 136.583~s and candidate pairs to 20,256.7, but its Recall is lower (0.8819 versus 0.9310); Precision rises from 0.0659 to 0.2328.}

\hl{On NCVR, Entity Eviction increases mean matching time (224.307~s versus 127.758~s) while candidate-pair count, $TP$, and Recall remain close to the baseline. Block Eviction reduces matching time to 87.459~s and candidate pairs to 107,462.3; observed $TP$ and Recall are unchanged at 38,662 and 0.9454, while Precision increases from 0.2298 to 0.3598.}

\hl{On MusicBrainz, Entity Eviction increases mean matching time to 1,636.732~s, with candidate-pair count and quality metrics close to the baseline. Block Eviction reduces matching time to 521.717~s and candidate pairs to 29,042.6. The observed $TP$ and Recall remain 20,578 and 0.1064, while Precision increases from 0.4814 to 0.7085. These results describe the evaluated datasets and settings; they do not show that the removed blocks were semantically irrelevant or that the same trade-offs generalize to other workloads.}

\hl{Table~\ref{tab:welch-t-tests-exp23} reports the exploratory two-sided Welch tests against the baseline. Since several metric samples are constant, the corresponding test statistics are undefined and shown as n/a. The p-values are not adjusted for multiple comparisons.}

\begin{table*}[!t]
\centering
\caption{\hl{V2 results for Experiments 1--3 (mean $\pm$ 95\% confidence-interval half-width across 10 replicas).}}
\label{tab:resultados_consolidados_rotacionada}
\scriptsize
\setlength{\tabcolsep}{4pt}
\renewcommand{\arraystretch}{0.95}
\begin{tabular}{llrrrrrrrr}
\hline
\hl{Dataset} & \hl{Method} & \hl{Blocking (s)} & \hl{Matching (s)} & \hl{Candidate pairs} & \hl{$TP$} & \hl{$FP$} & \hl{$FN$} & \hl{Recall} & \hl{Precision} \\
\hline
\hl{SCD} & \hl{Baseline} & \hl{$5.743\pm0.098$} & \hl{$155.856\pm1.014$} & \hl{$75{,}527.1\pm746.9$} & \hl{$4{,}978.2\pm3.4$} & \hl{$70{,}548.9\pm747.4$} & \hl{$368.8\pm3.4$} & \hl{$0.9310\pm0.0006$} & \hl{$0.0659\pm0.0007$} \\
\hl{SCD} & \hl{Entity Eviction} & \hl{$5.786\pm0.009$} & \hl{$216.571\pm1.417$} & \hl{$74{,}596.4\pm739.0$} & \hl{$4{,}972.9\pm3.4$} & \hl{$69{,}623.5\pm739.8$} & \hl{$374.1\pm3.4$} & \hl{$0.9300\pm0.0006$} & \hl{$0.0667\pm0.0007$} \\
\hl{SCD} & \hl{Block Eviction} & \hl{$5.933\pm0.057$} & \hl{$136.583\pm0.977$} & \hl{$20{,}256.7\pm183.1$} & \hl{$4{,}715.3\pm10.5$} & \hl{$15{,}541.4\pm186.4$} & \hl{$631.7\pm10.5$} & \hl{$0.8819\pm0.0020$} & \hl{$0.2328\pm0.0023$} \\
\hline
\hl{NCVR} & \hl{Baseline} & \hl{$31.064\pm0.079$} & \hl{$127.758\pm0.754$} & \hl{$168{,}228.1\pm801.4$} & \hl{$38{,}662.0\pm0.0$} & \hl{$129{,}566.1\pm801.4$} & \hl{$2{,}231.0\pm0.0$} & \hl{$0.9454\pm0.0000$} & \hl{$0.2298\pm0.0011$} \\
\hl{NCVR} & \hl{Entity Eviction} & \hl{$33.076\pm0.023$} & \hl{$224.307\pm1.166$} & \hl{$168{,}033.8\pm797.3$} & \hl{$38{,}662.0\pm0.0$} & \hl{$129{,}371.8\pm797.3$} & \hl{$2{,}231.0\pm0.0$} & \hl{$0.9454\pm0.0000$} & \hl{$0.2301\pm0.0011$} \\
\hl{NCVR} & \hl{Block Eviction} & \hl{$33.397\pm0.478$} & \hl{$87.459\pm1.869$} & \hl{$107{,}462.3\pm419.5$} & \hl{$38{,}662.0\pm0.0$} & \hl{$68{,}800.3\pm419.5$} & \hl{$2{,}231.0\pm0.0$} & \hl{$0.9454\pm0.0000$} & \hl{$0.3598\pm0.0014$} \\
\hline
\hl{MusicBrainz} & \hl{Baseline} & \hl{$360.866\pm0.494$} & \hl{$809.549\pm1.480$} & \hl{$42{,}744.2\pm98.5$} & \hl{$20{,}578.0\pm0.0$} & \hl{$22{,}166.2\pm98.5$} & \hl{$172{,}893.0\pm0.0$} & \hl{$0.1064\pm0.0000$} & \hl{$0.4814\pm0.0011$} \\
\hl{MusicBrainz} & \hl{Entity Eviction} & \hl{$378.246\pm5.188$} & \hl{$1{,}636.732\pm50.004$} & \hl{$42{,}658.1\pm102.3$} & \hl{$20{,}578.0\pm0.0$} & \hl{$22{,}080.1\pm102.3$} & \hl{$172{,}893.0\pm0.0$} & \hl{$0.1064\pm0.0000$} & \hl{$0.4824\pm0.0012$} \\
\hl{MusicBrainz} & \hl{Block Eviction} & \hl{$362.071\pm1.277$} & \hl{$521.717\pm2.928$} & \hl{$29{,}042.6\pm47.0$} & \hl{$20{,}578.0\pm0.0$} & \hl{$8{,}464.6\pm47.0$} & \hl{$172{,}893.0\pm0.0$} & \hl{$0.1064\pm0.0000$} & \hl{$0.7085\pm0.0011$} \\
\hline
\end{tabular}%
\end{table*}

\begin{table*}[!t]
\centering
\caption{\hl{Welch $t$-test $p$-values for Experiments 2 and 3 versus the V2 baseline.}}
\label{tab:welch-t-tests-exp23}
\scriptsize
\setlength{\tabcolsep}{5pt}
\begin{tabular}{llcc}
\hline
\hl{Dataset} & \hl{Metric} & \hl{Entity Eviction} & \hl{Block Eviction} \\
\hline
\hl{SCD} & \hl{Blocking time} & \hl{0.3515} & \hl{0.0019} \\
\hl{SCD} & \hl{Matching time} & \hl{$1.77\times10^{-22}$} & \hl{$4.75\times10^{-17}$} \\
\hl{SCD} & \hl{Candidate pairs} & \hl{0.0604} & \hl{$1.45\times10^{-18}$} \\
\hl{SCD} & \hl{$TP$} & \hl{0.0229} & \hl{$1.44\times10^{-14}$} \\
\hl{SCD} & \hl{$FP$} & \hl{0.0619} & \hl{$1.36\times10^{-18}$} \\
\hl{SCD} & \hl{Recall} & \hl{0.0229} & \hl{$1.44\times10^{-14}$} \\
\hl{SCD} & \hl{Precision} & \hl{0.0897} & \hl{$5.22\times10^{-19}$} \\
\hline
\hl{NCVR} & \hl{Blocking time} & \hl{$2.58\times10^{-14}$} & \hl{$1.12\times10^{-6}$} \\
\hl{NCVR} & \hl{Matching time} & \hl{$4.02\times10^{-26}$} & \hl{$1.22\times10^{-14}$} \\
\hl{NCVR} & \hl{Candidate pairs} & \hl{0.7020} & \hl{$2.39\times10^{-23}$} \\
\hl{NCVR} & \hl{$TP$} & \hl{n/a} & \hl{n/a} \\
\hl{NCVR} & \hl{$FP$} & \hl{0.7020} & \hl{$2.39\times10^{-23}$} \\
\hl{NCVR} & \hl{Recall} & \hl{n/a} & \hl{n/a} \\
\hl{NCVR} & \hl{Precision} & \hl{0.7016} & \hl{$1.00\times10^{-28}$} \\
\hline
\hl{MusicBrainz} & \hl{Blocking time} & \hl{$3.20\times10^{-5}$} & \hl{0.0705} \\
\hl{MusicBrainz} & \hl{Matching time} & \hl{$3.35\times10^{-11}$} & \hl{$1.67\times10^{-24}$} \\
\hl{MusicBrainz} & \hl{Candidate pairs} & \hl{0.1870} & \hl{$7.71\times10^{-44}$} \\
\hl{MusicBrainz} & \hl{$TP$} & \hl{n/a} & \hl{n/a} \\
\hl{MusicBrainz} & \hl{$FP$} & \hl{0.1870} & \hl{$7.15\times10^{-26}$} \\
\hl{MusicBrainz} & \hl{Recall} & \hl{n/a} & \hl{n/a} \\
\hl{MusicBrainz} & \hl{Precision} & \hl{0.1875} & \hl{$2.89\times10^{-35}$} \\
\hline
\end{tabular}

\hl{The tests use $\alpha=0.05$ and are not corrected for multiple comparisons. n/a indicates constant sample values for which the test statistic is undefined.}
\end{table*}




\subsection{Experiment 3: Evaluating the Global Eviction Technique}
\hl{The V2 files named ``descarte velho'' are the new outputs designated for the Global Eviction run. Their raw means are listed in Table~\ref{tab:resultados_consolidados_global_rotacionada}; however, the implementation does not provide an unbiased comparison of matching quality. During blocking, it uses the ground-truth mapping to select the Dataset 2 indices passed to the matcher for each Dataset 1 record. Because the list of correspondence lists is nonempty, the matcher's fallback scan over Dataset 2 is not used. Thus, the queries evaluated by this run are conditioned on known matches, rather than being processed as an independent stream in row order.}

\hl{The code also increments the reported $FP$ only when a predicted Dataset 1 identifier has no ground-truth entry at all. If that identifier has a ground-truth match to a different Dataset 2 identifier, the emitted incorrect pair is not counted as an $FP$. Accordingly, the raw candidate-pair, $TP$, and reported $FP$ counts do not form a valid full-dataset confusion matrix, and the reported Recall and Precision cannot be interpreted as unbiased matching-quality estimates. No inferential comparison of these quality metrics is reported here.}

\hl{For MusicBrainz, the raw mean output is 202,118.1 candidate pairs and 193,464.2 $TP$, against a denominator of 193,471 ground-truth rows used by the script. It also reports 4,420.6 $FP$, but candidate pairs minus $TP$ minus reported $FP$ equals 4,233.3, exposing the false-positive undercount. The near-denominator $TP$ result is consistent with the ground-truth-conditioned query selection and is not evidence that Global Eviction recovered almost all matches from an independent stream. The same protocol issue applies to DBLP and NCVR.}

\hl{The blocking and matching times and candidate counts in Table~\ref{tab:resultados_consolidados_global_rotacionada} are retained as raw execution observations, not as evidence of improved end-to-end performance: the matching workload differs from the baseline because ground truth determines which Dataset 2 records are queried. A valid quality comparison requires rerunning this policy with Dataset 2 processed independently of ground truth and counting every incorrect emitted pair as an $FP$.}


\begin{table*}[!t]
\centering
\caption{\hl{Raw V2 output for the Global Eviction run named ``descarte velho'' (mean $\pm$ 95\% confidence-interval half-width across 10 replicas).}}
\label{tab:resultados_consolidados_global_rotacionada}
\scriptsize
\setlength{\tabcolsep}{6pt}
\renewcommand{\arraystretch}{1.0}
\begin{tabular}{lrrrrrr}
\hline
\hl{Dataset} & \hl{Blocking (s)} & \hl{Matching (s)} & \hl{Candidate pairs} & \hl{$TP$ (reported)} & \hl{$FP$ (reported)} & \hl{$|MP|$ denominator} \\
\hline
\hl{SCD} & \hl{$6.313\pm0.103$} & \hl{$169.668\pm2.575$} & \hl{$24{,}625.6\pm125.5$} & \hl{$5{,}286.5\pm5.0$} & \hl{$839.5\pm29.8$} & \hl{5,347} \\
\hl{NCVR} & \hl{$52.412\pm1.305$} & \hl{$80.363\pm1.642$} & \hl{$128{,}081.4\pm599.4$} & \hl{$40{,}892.0\pm0.0$} & \hl{$43{,}857.3\pm322.0$} & \hl{40,893} \\
\hl{MusicBrainz} & \hl{$404.297\pm4.768$} & \hl{$541.225\pm4.598$} & \hl{$202{,}118.1\pm57.8$} & \hl{$193{,}464.2\pm1.0$} & \hl{$4{,}420.6\pm25.8$} & \hl{193,471} \\
\hline
\end{tabular}%
\end{table*}


\subsection{Experiment 4: Entity Lifetime in Blocks}

In order to identify the influence of the proposed eviction techniques (Entity Eviction and Block Eviction), the lifetime of each entity within the blocks was analyzed. For this purpose, after assigning entities to blocks, the time at which each entity was inserted into a block was compared with the current processing time of incoming entities. The results are presented in Table \ref{tab:comparacao_memoria_consolidada}. The shorter the lifetime, the faster entities are refreshed within the blocks, freeing space for the insertion of new entities.

Global Eviction was not evaluated in this experiment. Although entities are still assigned to blocks, eviction decisions are made independently of the block structure dynamics, based on global temporal criteria such as their lifetime in memory. Unlike Entity Eviction and Block Eviction, which are driven by block occupancy and saturation, Global Eviction removes entities simultaneously from all blocks in which they appear.

\begin{table}[h]
\centering
\caption{Comparison of memory lifetime (increments) for the SCD, NCVR, and MusicBrainz datasets.}
\label{tab:comparacao_memoria_consolidada}
\begin{tabular}{lccc}
\hline
\textbf{Method} & \textbf{SCD} & \textbf{NCVR} & \textbf{MusicBrainz} \\
\hline
Baseline & 876.91 & 16,555.14 & 20,858.32 \\
Entity Eviction & 814.19 & 14,031.15 & 19,253.74 \\
Block Eviction & 794.72 & 5,323.19 & 13,260.22 \\
\hline
\end{tabular}
\end{table}

The application of the Entity Eviction technique resulted in a reduction in the lifetime of entities within the blocks, with a decrease of 62 increments for the SCD dataset and approximately 2,500 increments for the NCVR dataset. These results demonstrate that the technique is effective in reducing the historical retention of entities, particularly in scenarios with higher block density and overload.

According to the presented results, the average lifetime of entities varies significantly across datasets. In the baseline, the SCD dataset shows approximately 880 increments, while NCVR reaches around 16,500 increments. This difference is influenced by the number of entities processed in each dataset: SCD contains approximately 2,500 entities, whereas NCVR includes around 40,000. The Entity Eviction technique effectively reduces lifetime by systematically removing outdated entities, preventing blocks from growing indefinitely. The magnitude of this reduction is directly proportional to the data volume: the larger the number of entities, the greater the impact of the eviction mechanism.

Regarding the Block Eviction technique, an even more significant reduction in lifetime was observed, especially for the NCVR dataset. Compared to the baseline, the reduction was 82.19 increments (9.37\%) for SCD and 11,231.95 increments (67.85\%) for NCVR. This indicates that by removing entire blocks, the technique acts more aggressively in memory management, freeing space quickly and efficiently. This approach is particularly advantageous in \textit{streaming} scenarios with a high rate of incoming entities, where data quickly become obsolete and the system must prioritize recent information to maintain both processing quality and speed \cite{ARAUJO2025102506}.

When analyzing the MusicBrainz dataset, a behavior similar to the other datasets is observed, but with significantly higher absolute values. While the baseline method shows an average lifetime of approximately 20,858 increments, Entity Eviction reduces this value to around 19,253, representing a moderate decrease consistent with the intermediate size of this dataset. In contrast, the Block Eviction technique again proves to be the most effective, reducing lifetime to approximately 13,260 increments—a reduction of more than 7,500 increments compared to the baseline. As with NCVR, this result indicates that block-level eviction is substantially more efficient for datasets with a large number of textual attributes and high combinatorial generation during the blocking process. This technique enables faster block refreshment and prevents prolonged retention, which is crucial for multimedia datasets such as MusicBrainz, where older records quickly become irrelevant for incremental matching tasks.


\subsection{Experiment 5: Evaluation with Progressive Offsets}

In order to evaluate the impact of using different offset values on the performance of the proposed techniques, experiments were conducted on the SCD, NCVR, and MusicBrainz datasets using a progressive exponential partitioning approach. The offsets were defined based on a geometric progression, varying across all datasets. The experiments were carried out using the Block Eviction technique, as it yielded the best overall results. This experiment analyzes how variations in the size of processed data increments affect both the quality of the ER task and computational efficiency. The use of geometrically increasing offsets allows the evaluation of different scenarios typical of web and streaming applications, where data volume and heterogeneity evolve over time \cite{ren2021online,hassani2019overview}.

An example of the application of such a technique can be found in \cite{webmedia}, which analyzed the impact of the pandemic on the usage of communication platforms, focusing on Discord\footnote{https://discord.com/}. During this period, a significant increase in interaction volume was observed, particularly during late-night hours (between 3 a.m. and 9 a.m.), leading to overload in data processing systems. In this context, applying different offsets could mitigate this issue by dynamically adjusting the volume of processed data per batch according to workload variations over time.

For the SCD dataset, a positive evolution in both Recall and Precision metrics was observed as the offset values increased. With offsetA = 50 and offsetB = 1,000, the obtained values were approximately 0.55 for Recall and 0.80 for Precision. When the offsets were increased (offsetA = 1,350 and offsetB = 27,000), Recall rose to 0.72 (a gain of 30.9\%), while Precision reached 0.82 (a gain of 2.5\%). These results indicate a significant improvement in coverage for the ER task, albeit with a more modest gain in Precision.

For the NCVR dataset, the experiments revealed a different behavior. Despite the substantial increase in the number of compared pairs (over 110 thousand pairs), the Recall value approached 1.0 for larger offsets, such as 9,000 and 27,000, indicating that the approach is capable of identifying nearly all true matches. However, Precision remained more modest, ranging between 0.41 and 0.59, suggesting that many incorrect entity pairs were included in the comparisons, a direct consequence of the expanded search space.

For the MusicBrainz dataset, the results obtained with different offset values reveal a stable behavior in terms of both quality and processing time. It can be observed that, even with progressively increasing offsets, both blocking and matching times remained nearly constant, varying only by a few seconds across experiments. Regarding the quality of generated pairs, both Recall and Precision exhibited consistent values across all scenarios, with Recall around 0.106 and Precision close to 0.83. These results indicate that, for this dataset, increasing offsets does not improve coverage (i.e., it does not increase the number of true positives), nor does it negatively impact Precision.

Tables \ref{tab:scd-results-transposed}, \ref{tab:ncvr-results-transposed}, and \ref{tab:music-brainz-results-transposed} present the relationship between offset values and the performance metrics obtained for the SCD, NCVR, and MusicBrainz datasets, respectively. For the SCD dataset, there is a tendency toward a balance between Recall and Precision, whereas for NCVR, Recall reaches values close to 1 for larger offsets, while Precision remains considerably lower, despite a slight increase. These results indicate that the choice of an appropriate offset should consider the trade-off between Recall and Precision, being strongly influenced by the structural characteristics of each dataset. For MusicBrainz, the results remain stable. In contexts with high heterogeneity or greater ambiguity among entities, such as NCVR, increasing the offset may benefit Recall, albeit with a potential decrease in Precision.

\begin{table}[h]
\centering
\caption{Experimental results for the SCD dataset with different offsets.}
\label{tab:scd-results-transposed}
\begin{tabular}{c|cccc}
\hline
\textbf{Metric} & \textbf{Exp. 1} & \textbf{Exp. 2} & \textbf{Exp. 3} & \textbf{Exp. 4} \\
\hline
Offset A & 50 & 150 & 450 & 1350 \\
Offset B & 1000 & 3000 & 9000 & 27000 \\
Blocking (s) & 4.992 & 4.986 & 4.986 & 4.992 \\
Matching (s) & 112.374 & 113.958 & 117.984 & 128.790 \\
Recall & 0.5504 & 0.5633 & 0.6053 & 0.7226 \\
Precision & 0.8008 & 0.8037 & 0.7952 & 0.8249 \\
\hline
\end{tabular}
\end{table}

\begin{table}[h]
\centering
\caption{Experimental results for the NCVR dataset with different offsets.}
\label{tab:ncvr-results-transposed}
\begin{tabular}{c|cccc}
\hline
\textbf{Metric} & \textbf{Exp. 1} & \textbf{Exp. 2} & \textbf{Exp. 3} & \textbf{Exp. 4} \\
\hline
Offset A & 1000 & 3000 & 9000 & 27000 \\
Offset B & 1000 & 3000 & 9000 & 27000 \\
Blocking (s) & 32.610 & 32.214 & 32.082 & 32.028 \\
Matching (s) & 86.124 & 85.938 & 85.896 & 88.044 \\
Recall & 0.8288 & 0.8533 & 1.0000 & 1.0000 \\
Precision & 0.4122 & 0.3670 & 0.5550 & 0.5973 \\
\hline
\end{tabular}
\end{table}

\begin{table}[h]
\centering
\caption{Experimental results for the MusicBrainz dataset with different offsets.}
\label{tab:music-brainz-results-transposed}
\begin{tabular}{c|cccc}
\hline
\textbf{Metric} & \textbf{Exp. 1} & \textbf{Exp. 2} & \textbf{Exp. 3} & \textbf{Exp. 4} \\
\hline
Offset A & 5000 & 15000 & 45000 & 135000 \\
Offset B & 5000 & 15000 & 45000 & 135000 \\
Blocking (s) & 357.67 & 360.52 & 358.58 & 363.38 \\
Matching (s) & 776.35 & 782.71 & 775.16 & 794.58 \\
Recall & 0.1064 & 0.1064 & 0.1064 & 0.1064 \\
Precision & 0.8349 & 0.8320 & 0.8325 & 0.8321 \\
\hline
\end{tabular}
\end{table}

\subsection{Experiment 6: Impact of the Absence of Eviction Techniques} \label{exp6}

This experiment aims to evaluate the impact of the absence of eviction techniques on the ER task, enabling a direct comparison with the scenario in which a single eviction of the oldest entity is applied (baseline). In the no-eviction scenario, all entities remain active throughout the entire processing, leading to continuous growth of the search space and an increase in the number of comparisons. The results of this evaluation are presented in Table \ref{tab:baseline_vs_sem_descarte}.

For the SCD dataset, as it is the smallest dataset evaluated, the absence of eviction does not produce a significant impact on computational cost. The blocking time is 8.4~s and the matching time is 156.5~s, with 71,256 generated pairs. These values are close to those observed in the baseline (10.56~s for blocking, 171.96~s for matching, and 71,883.5 pairs). Similarly, the quality metrics remain virtually unchanged, with Recall of 92\% and Precision of 66\% in the no-eviction scenario, compared to 93\% and 66\% in the baseline. This behavior indicates that, for small datasets, eviction has a limited effect on both efficiency and quality.

For the NCVR dataset, characterized by a larger volume and greater diversity of records, the absence of eviction leads to a significant increase in computational cost. The matching time increases from 137.04~s in the baseline to 529.1~s in the no-eviction scenario, while the number of pairs rises from 191,632.1 to 289,125. Although Recall increases from 85.5\% to 94.5\%, a reduction in Precision is observed, decreasing from 30.5\% to 23.6\%. This result indicates that the gain in Recall comes at the expense of a large number of incorrect comparisons, making the process less selective and more computationally expensive.

The most critical behavior is observed for the MusicBrainz dataset. Without eviction, the matching time reaches approximately 30,031.2~s, with 442,586 generated pairs, whereas in the baseline these values are 1,229.53~s and 41,781 pairs, respectively. Moreover, Precision drops sharply, from 65.7\% in the baseline to only 8\% in the no-eviction scenario, while Recall remains virtually unchanged (10.6\% vs. 10.6\%). These results demonstrate that, for larger and more heterogeneous datasets, the absence of eviction makes the process impractical, both in terms of computational cost and the quality of the produced matches.

Overall, the comparison between the baseline and the no-eviction scenario demonstrates that eviction techniques are essential for controlling the growth of the search space, reducing unnecessary comparisons, and preserving result quality. For small datasets, their impact is limited; however, as data volume and complexity increase, eviction becomes a critical component to ensure scalability, stability, and adequate performance of the ER task.

\begin{table}[!t]
\centering
\caption{Comparison between the baseline and the no-eviction scenario.}
\label{tab:baseline_vs_sem_descarte}

\scriptsize
\setlength{\tabcolsep}{3pt}
\renewcommand{\arraystretch}{0.95}
\resizebox{0.48\textwidth}{!}{%
\begin{tabular}{llcccccc}
\hline
\textbf{Dataset} & \textbf{Scenario} & \textbf{Block (s)} & \textbf{Match (s)} & \textbf{\#Pairs} & \textbf{TP} & \textbf{Recall} & \textbf{Precision} \\
\hline

\multirow{2}{*}{SCD}
& Baseline      & 10.56 & 171.96  & 71883.5 & 5977.3 & 0.9300 & 0.6672 \\
& No Eviction    & 8.40  & 156.50  & 71256.0 & 4971.0 & 0.9297 & 0.6661 \\
\hline

\multirow{2}{*}{NCVR}
& Baseline      & 33.36 & 137.04  & 191632.1 & 34980.6 & 0.8550 & 0.3050 \\
& No Eviction    & 29.10 & 529.10  & 289125.0 & 38662.0 & 0.9454 & 0.2361 \\
\hline

\multirow{2}{*}{MusicBrainz}
& Baseline      & 346.06 & 1229.53 & 41781.0  & 20578.0 & 0.1060 & 0.6570 \\
& No Eviction    & 577.80 & 30031.20 & 442586.0 & 20578.0 & 0.1064 & 0.0885 \\
\hline

\end{tabular}%
}
\end{table}

\section{Conclusions and Future Work} \label{sec:conclusao}

\hl{The V2 evaluation compares the AB baseline with time-aware Entity Eviction and Block Eviction on SCD, NCVR, and MusicBrainz. The results answer the research question in a dataset-dependent manner: the policies change processing time, candidate counts, Recall, and Precision differently, and the measurements do not support a general claim that matching quality is preserved.}

\hl{In the V2 comparisons, Entity Eviction increased mean matching time relative to the baseline on all three datasets. Candidate-pair counts and Precision changed only slightly, while SCD showed a small reduction in $TP$ and Recall and NCVR and MusicBrainz retained their observed $TP$ and Recall. These results do not demonstrate a general runtime or memory-stability benefit.}

\hl{The small quality-metric changes for Entity Eviction should be interpreted as observations for the tested datasets and configuration, not as evidence that the policy preserves quality in other settings. The reported experiment does not directly measure memory use or establish operational stability.}

\hl{Block Eviction reduced mean matching time and candidate-pair counts relative to the baseline on all three datasets, and increased Precision. On SCD, Recall decreased from 0.9310 to 0.8819; on NCVR and MusicBrainz, the observed $TP$ and Recall were unchanged. The SCD result is a measured quality trade-off, not evidence that block removal preserves all relevant matches.}

\hl{The increased Precision is consistent with a more selective set of emitted candidate pairs under the implemented matcher. The experiments do not establish why particular blocks contain false positives, nor do they establish performance in large-scale or distributed deployments beyond the tested single-machine settings.}

\hl{The new V2 run named ``descarte velho'' cannot support a matching-quality conclusion: its script uses ground truth to select Dataset 2 queries and undercounts false-positive pairs. Its reported Recall and Precision, including the near-one MusicBrainz Recall, are therefore excluded from claims about Global Eviction. A quality comparison requires a rerun with query selection independent of ground truth and pair-wise $FP$ counting.}

\hl{The separate entity-lifetime measurements report shorter average lifetimes under Entity Eviction and Block Eviction than under the baseline; they do not demonstrate that shorter lifetime caused the observed runtime or quality differences. For the evaluated workloads, Block Eviction reduced matching time and candidate counts, while its Recall effect depended on dataset. The current evidence does not establish a universally preferred policy.}

\hl{These findings are limited to the evaluated datasets, parameter settings, and single-machine environment. The Global Eviction run named ``descarte velho'' requires a corrected, ground-truth-independent evaluation before its quality can be compared with the baseline. Further work should also evaluate parameter sensitivity and distributed or production streaming settings.}

\begin{enumerate}
    \item \textbf{Adaptive hybrid techniques.}  
    Investigate the dynamic combination of Global Eviction, Entity Eviction, and Block Eviction, enabling the selection of the most appropriate strategy according to data stream characteristics;

    \item \textbf{Semantic and statistical criteria.}  
    Incorporate criteria based on semantic and statistical relevance to avoid removing important entities and improve matching quality;

    \item \textbf{Distributed environments.}  
    Evaluate the impact of eviction techniques in distributed systems, considering aspects such as scalability, memory usage, and latency;

    \item \textbf{Machine learning.}  
    Use machine learning models to guide eviction decisions in an adaptive manner, based on the historical behavior of the data.
\end{enumerate}


\bibliographystyle{IEEEtran}
\bibliography{access-bibliography}


\begin{IEEEbiography}[{\includegraphics[width=1in,height=1.25in,clip,keepaspectratio]{igor-foto.jpg}}]{Igor de Sousa Pereira} received the M.S. degree in Computer Science from the Federal University of Campina Grande (UFCG), Brazil. He is currently a Researcher at the UFCG, where he conducts research at the Laboratory of Artificial Intelligence and Dedicated Architectures (LIAD) in collaboration with a multinational company. His research focuses on state-of-the-art predictive methods and intelligent data-driven solutions. His research interests include artificial intelligence, entity resolution, data analytics, machine learning, and data engineering.
\end{IEEEbiography}

\newpage

\begin{IEEEbiography}[{\includegraphics[width=1in,height=1.25in,clip,keepaspectratio]{tiago_brasileiro.jpeg}}]{Tiago Brasileiro Araújo} is a Professor at the Federal Institute of Paraíba (IFPB) and a researcher affiliated with Tampere University (TUNI), Finland, the VIRTUS Competence Center, and the Applied Intelligent Computing Laboratory (LACINA) at the Federal University of Campina Grande (UFCG), Brazil. He completed his PhD in Computer Science at UFCG, Brazil, followed by a second PhD at Tampere University, Finland. His research focuses on entity resolution, cloud computing, artificial intelligence, and data analytics, with a growing emphasis on the socio-technical aspects of data management, particularly fairness and transparency. He has co-authored papers in top-tier venues and actively collaborates on both research and industry-driven projects in AI-driven solutions, data analytics, and decision-support systems. %He has co-authored papers in top-tier venues and frequently contributes to discussions on data analytics and AI, especially in the areas of data matching and intelligent systems.
\end{IEEEbiography}

\begin{IEEEbiography}[{\includegraphics[width=1in,height=1.25in,clip,keepaspectratio]{Edu trabalho_IMG 7521-1 - Copia.jpg}}]{Carlos Eduardo S, Pires} is a Professor at the Federal University of Campina Grande (UFCG), Brazil. He received his Ph.D. in Computer Science from the Federal University of Pernambuco (UFPE), Brazil. He leads research, development, and innovation projects at the Applied Intelligent Computing Laboratory (LACINA) in collaboration with national and multinational companies. His research interests include Data Management, with an emphasis on Data Quality, Data Integration, Knowledge Discovery, and Big Data.
\end{IEEEbiography}

\EOD

\end{document}